% Back cover (always lands on an even page)
\cleardoublepage
\thispagestyle{empty}\null\clearpage
\thispagestyle{empty}
\begin{tikzpicture}[remember picture, overlay]
    \fill[coverPaper] (current page.south west) rectangle (current page.north east);
    \fill[coverInk] (current page.north west) rectangle ([xshift=2.1cm]current page.south west);
    \fill[coverGold] ([xshift=2.1cm]current page.north west) rectangle ([xshift=2.25cm]current page.south west);
    \node[anchor=north, text width=13.5cm, align=right, lawn=coverInk]
        at ([xshift=1.1cm, yshift=-3.2cm]current page.north) {%
        {\khatt\bfseries\LARGE\color{coverInk} تعلّم الخوارزميات بحلّ المسائل}\par
        \vspace{2pt}{\color{coverGold}\rule{5cm}{0.8pt}}\par\vspace{10pt}
        \color{coverInk}\large
        يرافق هذا الكتاب طالب السنة الأولى في مقياس الخوارزميات ولغة \en{C}،
        من أول متغير إلى الأشجار الثنائية، عبر \textbf{\en{650} تمرينًا محلولًا}
        مأخوذة من مسائل واقعية في \textbf{التجارة الإلكترونية} و\textbf{علم البيانات}.\par
        \vspace{10pt}
        \faCheckCircle\ \en{13} فصلًا، في كل فصل \en{50} تمرينًا: \en{10} سهلة، \en{30} متوسطة، \en{10} صعبة.\par
        \faCheckCircle\ لكل تمرين: نص المسألة، خطوات الحل، برنامج \en{C} مشروح سطرًا بسطر،
        نتيجة التنفيذ الحقيقية، ومخطط انسيابي.\par
        \faCheckCircle\ كل الخوارزميات مكتوبة يدويًا، دون أي دالة جاهزة.\par
        \faCheckCircle\ برامج مترجمة ومختبرة بالمعيار \en{C99}.\par
        \vspace{14pt}
        {\khatt\bfseries\Large\color{coverInk}\teacher}\par
        {\color{coverInk!70} أستاذ مقياس الخوارزميات وهياكل المعطيات}};
    % ISBN / barcode area
    \node[draw=coverInk!60, fill=white, minimum width=5.2cm, minimum height=2.8cm,
        anchor=south east, font=\small\color{coverInk!70}]
        at ([xshift=-1.6cm, yshift=1.8cm]current page.south east)
        {\en{ISBN}\quad\en{\bookISBN}};
\end{tikzpicture}
